% Plan: development/outline.md, section 6, item I. Internal script names
% may appear here, in neutral form.
% Sources: R/publication/reproduction/README.md and commands.json (commands,
% recorded wall times), R/publication/integration/runtime-evidence-r1.json,
% the family dossiers and critiques, development/reviews/*.md,
% development/eg-audit/out/{run_audit,compare}.log (per-instance sums of the
% audited eg run: eg_int_s 465 s, eg_disc_s 1,456 s, eg_disc2_s 14,782 s),
% R/publication/reproduction/water-audit/logs/, R/publication/reproduction/
% network/logs/ann.*.time, development/reviews/round1/kan-guard/ (guarded KAN
% replay) and development/reviews/round1/sol-numbers.md (rocket and GAMS/OSIL
% reruns). Expected outputs point to the generated tables and to
% data/numbers.json. Tiers follow the rule of Section 10 (recorded wall time
% per instance, summed over its parts). Round-1 revision (G7-12): status per
% register row, trust labels defined here, archive layout, setup sequence.
% The stored-model table is generated (tables/tab-sem-hashes.tex).
\section{Reproduction guide and claim register}\label{app:repro}

The files and scripts that replay each result are named here; the main text avoids these names.
As \cref{sec:semantics-protocol} states, AI agents working under the authors' direction wrote the implementations and ran the checks; that section defines ``separately written''.
Paths are relative to the root of the archive (\archiveDOI), which keeps the layout of our working tree, including working folder names such as \nolinkurl{open-instances-wave3/} or \nolinkurl{dossiers/}: \nolinkurl{$R} denotes its folder \nolinkurl{research-20260929/} and \nolinkurl{$P} its folder \nolinkurl{paper-open-minlplib/}.
\Cref{tab:repro-register} lists the replay entry points; times are recorded wall times on a shared machine, and \cref{sec:repro} defines the tiers by these times.

\subsection{Setup}\label{app:repro-setup}

The guide \nolinkurl{$P/artifact/README.md} describes the files of the artifact, the environment, the long replays and the earlier reproduction package in \nolinkurl{$R/publication/reproduction/}.
The checkers need Python~3.13.11 and the package versions pinned in \nolinkurl{$R/publication/reproduction/requirements.txt}; the \inst{dtoc5} checker also needs the GMP library, and the solver campaign needs GAMS~54.3.1 and solver licences.
Several scripts write into their own \nolinkurl{logs/} folders and some replace stored evidence, so every command runs in a disposable copy \nolinkurl{$WORK} of the archive, with at most four single-threaded processes at a time.
A session starts as follows; the unpacked archive \nolinkurl{$ARCHIVE} is only read.
\begin{verbatim}
ARCHIVE=/path/to/unpacked/archive          # read only
WORK=$(mktemp -d /tmp/minlp-work-XXXXXX)   # disposable: rm -rf "$WORK"
cp -r "$ARCHIVE"/. "$WORK"/ && cd "$WORK"
export HOME="$WORK/home" XDG_CACHE_HOME="$WORK/home/.cache"
export TMPDIR="$WORK/tmp" PIP_NO_CACHE_DIR=1 PYTHONDONTWRITEBYTECODE=1
R="$WORK/research-20260929"; P="$WORK/paper-open-minlplib"
export MINLP_REPO_ROOT="$WORK" EG_AUDIT_R="$R" MINLPLIB_OSIL_ROOT="$WORK/osil"
export OMP_NUM_THREADS=1 OPENBLAS_NUM_THREADS=1 MKL_NUM_THREADS=1
export NUMEXPR_NUM_THREADS=1
mkdir -p "$TMPDIR" "$MINLPLIB_OSIL_ROOT" "$HOME/.cache/minlplib/minlplib"
cp "$R"/publication/minlplib-status/pages/models/osil/*.osil \
   "$MINLPLIB_OSIL_ROOT"
ln -s "$MINLPLIB_OSIL_ROOT" "$HOME/.cache/minlplib/minlplib/osil"
python3 "$P/artifact/check_claims.py" --repo-root "$WORK"
python3 -m venv "$WORK/venv" && . "$WORK/venv/bin/activate"
python -m pip install -r "$R/publication/reproduction/requirements.txt"
python "$P/artifact/prepare_paper_inputs.py"
cd /tmp && python "$P/data/make_tables.py"
\end{verbatim}
The archive holds the 69 model files of the paper (\cref{tab:sem-hashes,tab:audit-inputs}) in the folder copied above.
The relocated review scripts read them from \nolinkurl{MINLPLIB_OSIL_ROOT}, the other research scripts from \nolinkurl{~/.cache/minlplib/minlplib/osil/}, which, with \nolinkurl{HOME} inside the copy, is a link to the same folder.
Nothing is then written outside \nolinkurl{$WORK} and the folders \nolinkurl{/tmp/minlp-artifact-*} of the short-check runner, except by two archived drivers: the driver \nolinkurl{run_replays.py} of the guarded KAN replay expects its tree in \nolinkurl{/tmp/kan-guard} and six cores, and the campaign driver names the GAMS installation of the original machine.
The README gives a replay of the guarded KAN searches inside \nolinkurl{$WORK}, one instance per process.
The validator runs first because \nolinkurl{make_tables.py} rewrites the number file and \nolinkurl{check.log} with the date of the run.
\nolinkurl{prepare_paper_inputs.py} runs the archived \nolinkurl{prepare_inputs.py}, which restores the inputs that version control ignores (listed points and archived instance pages) and checks the SHA-256 of every input and model; it exits with status~0 if the 69 models are present with their recorded hashes and the only missing models are the 26 others of the manifest, none used by the paper, and it never downloads.
\nolinkurl{make_tables.py} (standard library only, about 5\,s) checks every bound, primal, gap and derived margin display that it generates against its exact value in rational arithmetic, writes the number file \nolinkurl{$P/data/numbers.json} and \nolinkurl{$P/data/check.log}, which lists every check and the SHA-256 of every file read, and exits with status~1 if a check fails (\cref{app:displays}).
This recipe and the short checks below were run as written in a fresh copy (log \nolinkurl{$P/artifact/logs/recipe-test.log}); the long replays of \cref{tab:repro-register} were not repeated under it.

\paragraph{Trust labels.}
The last column of \cref{tab:repro-register} names the components of the trust base (\cref{sec:semantics-trust}) that a result needs:
\trust{int}, Python integer and rational arithmetic and integer square roots;
\trust{fp}, \cref{hyp:fp-ieee};
\trust{mp}, mpmath~1.3.0 interval arithmetic for the named primitives;
\trust{read}, the reading of \cref{def:sem-model} by the family's readers (\cref{sec:semantics-protocol});
\trust{code}, the correctness of the named checking programs (inspected and rerun, not formally verified);
\trust{thm}, the classical theorems used (mean value, intermediate value, Brouwer and Sturm theorems, the Lindemann--Weierstrass theorem, series remainder bounds).
Every row uses \trust{read}, \trust{code} and \trust{thm}.

\paragraph{Claim index and short checks.}
The file \nolinkurl{$P/artifact/claims.json}, the machine-readable form of the claim register, has 65 entries, one for each of the 27 rows of \cref{tab:repro-register} and one for each of the 38 reported values of \cref{tab:claims,tab:claims-campaign}, with the displayed values, the SHA-256 of every input, code and output file and the status of each instance (the README lists all fields).
The validator \nolinkurl{$P/artifact/check_claims.py} uses only the Python standard library; it recomputes every hash and checks the paths, the labels, the displayed values and status labels of every instance, and that \cref{tab:repro-register} and \nolinkurl{RUNS.md} list the same rows.
On the final sources of this version it reported 5,632 SHA-256 references to 1,918 distinct files, all valid (log \nolinkurl{$P/artifact/logs/check_claims.log}); it fails on any hash that is out of date.
It validates the sources and the evidence (hashes of inputs, code and outputs, paths and labels); it does not identify the PDF files, whose SHA-256 hashes \nolinkurl{$P/artifact/RELEASE.md} records.
The runner \nolinkurl{$P/artifact/run_short_checks.py} copies what it needs into a new folder under \nolinkurl{/tmp} and runs 15 short checks there (\inst{lnts}, \inst{dtoc5}, the three \texttt{r3} KAN searches, the \inst{eg} auditor tests and review scripts); all 15 pass, with every stored field reproduced except elapsed time and CPU affinity (log \nolinkurl{$P/artifact/logs/short-checks.json}), and \nolinkurl{--eg-int} adds the regeneration and audit of the \inst{eg_int_s} tree (324\,s; 33,385 leaves; no violation).
\nolinkurl{$P/artifact/path-edits.json} lists the edits that made the review scripts, the audit and \nolinkurl{gibbs_cert.py} read their input roots from the variables above; none changes a numerical formula.

\subsection{Stored models}\label{app:repro-models}

The stored models of \cref{def:sem-model} are the OSIL files of \cref{tab:sem-hashes}; the archive records the full hashes, and \nolinkurl{prepare_inputs.py} rejects any file whose hash differs.

% Generated by paper-open-minlplib/data/make_tables.py from data/numbers.json.
% Do not edit by hand; rerun the script. Requires booktabs, array and rotating;
% \inst{} typesets an instance name; underscores are not escaped (macros.tex).
\begin{table}[htbp]
\centering
\footnotesize
\setlength{\tabcolsep}{3pt}
\caption{The 43 stored OSIL models: $n$ ($n_{\mathrm{int}}$)/$m$, variables (integer variables)/rows, and the first 16 hexadecimal digits of the SHA-256 hash. The files were fetched on 2026-09-29 and found unchanged on 2026-10-02.}
\label{tab:sem-hashes}
\begin{tabular}{@{}lll@{\qquad}lll@{}}
\toprule
instance & $n$ ($n_{\mathrm{int}}$)/$m$ & SHA-256 prefix & instance & $n$ ($n_{\mathrm{int}}$)/$m$ & SHA-256 prefix \\
\midrule
\inst{lnts50} & 256/200 & \texttt{3d8234b338afc4cd} & \inst{etamac} & 97/70 & \texttt{6716cc88e1c4c61f} \\
\inst{lnts100} & 506/400 & \texttt{284bddc374db070c} & \inst{pindyck} & 116/96 & \texttt{c7b3b5e4a2c46b3e} \\
\inst{lnts200} & 1006/800 & \texttt{c2dde9ef8ede3c16} & \inst{powerflow0030p} & 236/555 & \texttt{5c4b386b1f9a721a} \\
\inst{lnts400} & 2006/1600 & \texttt{96eb445643601734} & \inst{powerflow0039p} & 282/657 & \texttt{0a27073f39f6cc49} \\
\inst{dtoc5} & 99999/49999 & \texttt{ac7b2d9472dbbfe9} & \inst{powerflow0039r} & 282/473 & \texttt{d8ba3132cb69744e} \\
\inst{optcdeg2} & 150002/100000 & \texttt{6a9bcd91abdbe1b2} & \inst{hvycrash} & 201/150 & \texttt{27a0a05abecc0d64} \\
\inst{lukvle10} & 1000/998 & \texttt{1eb1236df017099f} & \inst{eg_int_s} & 8 (3)/28 & \texttt{a462a01b98801f45} \\
\inst{chain50} & 102/51 & \texttt{f6b2b409aa4f3d00} & \inst{eg_disc_s} & 8 (4)/28 & \texttt{5f21b8d743b12c22} \\
\inst{chain100} & 202/101 & \texttt{45cb302006643309} & \inst{eg_disc2_s} & 8 (3)/28 & \texttt{8849e06b478f5b93} \\
\inst{chain200} & 402/201 & \texttt{634acc6242ce83fa} & \inst{waterno2_06} & 996 (54)/1234 & \texttt{3b0b316d7a9fa834} \\
\inst{chain400} & 802/401 & \texttt{7cd93486e4610aa4} & \inst{waterno2_09} & 1494 (81)/1852 & \texttt{015932376b18ded5} \\
\inst{catmix100} & 303/200 & \texttt{9f5b5bae82032ecc} & \inst{waterno2_12} & 1992 (108)/2470 & \texttt{add485a67c22faee} \\
\inst{catmix200} & 603/400 & \texttt{9194135d049ce928} & \inst{waterno2_18} & 2988 (162)/3706 & \texttt{968ff496844df939} \\
\inst{catmix400} & 1203/800 & \texttt{3dfc4d7f2886fa12} & \inst{waterno2_24} & 3984 (216)/4942 & \texttt{91d2612f435cd751} \\
\inst{catmix800} & 2403/1600 & \texttt{40b572a4b80446a5} & \inst{ann_cumene_tanh} & 794/790 & \texttt{13a44f74762fdb47} \\
\inst{camshape100} & 199/200 & \texttt{9e17da72f0cfe538} & \inst{kan_r3_h1_n4} & 1129 (288)/1478 & \texttt{2991cd53706de81f} \\
\inst{camshape200} & 399/400 & \texttt{76d6141d111fb1a0} & \inst{kan_r3_h1_n5} & 1410 (360)/1847 & \texttt{203647dbe25965ee} \\
\inst{camshape400} & 799/800 & \texttt{a53174a46ad923ef} & \inst{kan_r3_h1_n9} & 2534 (648)/3323 & \texttt{267fc6596ea164a0} \\
\inst{camshape800} & 1599/1600 & \texttt{d0692d0f81211081} & \inst{kan_r5_h1_n3} & 838 (216)/1121 & \texttt{75d5998686158c0d} \\
\inst{ex6_2_5} & 9/3 & \texttt{eed480570b4fa8c1} & \inst{kan_r5_h1_n5} & 1392 (360)/1867 & \texttt{4288259a89d635d9} \\
\inst{ex6_2_7} & 9/3 & \texttt{c7f7cd3756d230cf} & \inst{kan_r5_h1_n8} & 2223 (576)/2986 & \texttt{595f66966d00affc} \\
\inst{pricing050} & 50/5 & \texttt{56d3e00fcac90c09} &  & &  \\
\bottomrule
\end{tabular}
\end{table}

% tab:trust-full is input here, before the register longtable, so that it prints
% before Table S38 (a float deferred past the start of a longtable prints after it).
% Generated by paper-open-minlplib/data/make_tables.py from data/numbers.json.
% Do not edit by hand; rerun the script. Requires booktabs, array and rotating;
% \inst{} typesets an instance name; underscores are not escaped (macros.tex).
\begin{sidewaystable}
\centering
\scriptsize
\setlength{\tabcolsep}{4pt}
\caption{Trust base, data readings and primal constructions by certificate family (the full form of \cref{tab:trust}).
Dual: tags and level (\cref{sec:semantics-trust}), trusted primitives, data reading (\cref{app:semantics-codes}; ``X / Y'': codes differ).
Status as in \cref{tab:trust}.
Read first: whether, according to the session records, the agent session that wrote the later of the two implementations had read the code of the earlier one before its own results (\emph{format only}: only to learn the file format; \emph{not checked}: records not examined).
Primal: construction (\cref{sec:points-constructions}); mpmath-free proofs: \cref{tab:points-all}.}
\label{tab:trust-full}
\begin{tabular}{@{}>{\raggedright\arraybackslash}p{2.7cm}>{\raggedright\arraybackslash}p{4.7cm}>{\raggedright\arraybackslash}p{5.0cm}>{\raggedright\arraybackslash}p{2.3cm}>{\raggedright\arraybackslash}p{1.6cm}>{\raggedright\arraybackslash}p{1.4cm}>{\raggedright\arraybackslash}p{4.9cm}@{}}
\toprule
family & dual: arithmetic, trusted primitives; data reading & displayed dual certified by; second implementation & status & read first & level & primal construction and arithmetic \\
\midrule
\inst{lnts50}--\allowbreak\inst{lnts400} & \arith{E} integers, rationals, integer square roots; exact & two integer-only codes certify the enclosure of $\vstar$ & verified & yes & hand+\allowbreak stored & attaining point (\cref{thm:lnts-opt}); Krawczyk points as a second witness \\
\inst{dtoc5} & \arith{E} rationals, sums in GMP; exact & two exact codes (full rational $q(\lambda)$; fractions rounded up to $2^{-256}$) certify the displays & verified & yes & hand+\allowbreak stored & explicit rational point \arith{E} \\
\inst{optcdeg2} & \arith{E} exact stage minimization (Sturm sequences); exact & one exact code; a separately written interval code certifies 293.8760750958728 & weaker second & yes & stored & intermediate-value bracket; integer intervals \arith{E} \\
\inst{camshape100}--\allowbreak\inst{camshape800} & \arith{E} exact rational checks; exact & two exact codes and one with directed rounding agree to all printed digits & verified & yes & hand+\allowbreak stored & envelope point \arith{E} \\
\inst{lukvle10} & \arith{I} exp, log (both codes); exact / outward & one code; a separately written one certifies 352.238025369202 & weaker second & no & rerun & triangular definitions from 640-digit seeds \\
\inst{chain50}--\allowbreak\inst{chain400} & \arith{I} sqrt, log (both codes); outward & two codes certify the same binary64 bound & verified & yes & rerun & explicit point in a real quadratic field \arith{E} \\
\inst{catmix100}--\allowbreak\inst{catmix800} & \arith{F} $+,-,\times,\div$ only; exact, then outward & one code; a separately written one certifies weaker bounds & weaker second & yes & rerun & exact rational states \arith{E} \\
\inst{ex6_2_5}, \inst{ex6_2_7} & \arith{I} log; exact, then outward & one code; a separately written one certifies slightly stronger bounds & verified & yes & rerun & rational point; objective in exact rational intervals \\
\inst{pricing050} & \arith{I} exp; outward & one code; a separately written one reproduces it & verified & yes & stored & interior point; row slacks \arith{I} (two codes) \\
\inst{etamac} & \arith{I} exp, log; exact / outward & one code; a separately written one certifies only $-15.294675643368096$ & weaker second & no & stored & triangular definitions \arith{I} (one code) \\
\inst{pindyck} & \arith{F}+\arith{I} concavity proof; bound \arith{E} from \arith{I} enclosures; exact / outward & concavity by two codes; bound from one code's \arith{I} enclosures, a slightly weaker one from the other's & weaker second & yes & stored & interval Newton per state \arith{I} (two codes) \\
the three \inst{powerflow} instances & \arith{E} Lagrangian; PSD proof by exact LDL$^\top$; exact & two exact codes & verified & 0039p, 0039r: yes; 0030p: format only & stored & Krawczyk test on an active-set system; dyadic re-proof \\
\inst{hvycrash} & written proof of an identity, checked in exact rationals; exact / outward & written proof; one code checks the identities exactly & proved & not checked & hand & backward triangular definitions; feasibility by a written proof; two witnesses checked in \arith{I} \\
the three \inst{eg} instances & \arith{F} under \cref{hyp:fp-ieee} with \cref{lem:eg-padding}; exp and powers checked by the accuracy auditor (complete run passed); rounded & leaf re-certifier on all leaves; second certificate for \inst{eg_int_s}, \inst{eg_disc_s}, part of \inst{eg_disc2_s} & verified; \inst{eg_disc2_s}: partial second & not checked & rerun & interior points; dyadic intervals \arith{E} \\
\inst{waterno2_06}--\allowbreak\inst{waterno2_24} & \arith{F} / \arith{E} rational node bounds; Python's correctly rounded \texttt{math.fsum}; outward / exact & either of two codes yields every value & verified & yes & rerun & exact algebraic point \arith{E} \\
\inst{ann_cumene_tanh} & \arith{F}, no library transcendental function; outward / rounded & two codes on one partition & proved; the second code copied an idiom of the first & read; order not recorded & rerun & forward triangular definitions; rational intervals \\
KAN (six instances) & \arith{F}; one rigorous exp and interval core shared by both codes; outward & reported $L$: the weaker of two paths; path~(II) assumes no NaN; path~(I) replayed with guarded bounds & verified; shared core & not checked & rerun & points of $\RP$; rational intervals \\
\bottomrule
\end{tabular}
\end{sidewaystable}


\paragraph{Expression nodes.}
Besides \texttt{variable}, \texttt{number} (all files but \inst{ann_cumene_tanh}), \texttt{sum}, \texttt{product} and \texttt{negate}, the nodes of \cref{def:sem-model} occur as follows:
\texttt{divide} in \inst{ex6_2_5}, \inst{pindyck}, \inst{hvycrash} and the KAN instances;
\texttt{square} in \inst{lukvle10}, \inst{chain}, \inst{pricing050}, \inst{powerflow0030p}, \inst{powerflow0039p} and \inst{eg};
\texttt{sqrt} in \inst{chain};
\texttt{power} in \inst{lukvle10}, \inst{etamac}, \inst{pindyck}, \inst{pricing050} and \inst{waterno2};
\texttt{exp} in \inst{pricing050}, \inst{eg} and the KAN instances;
\texttt{ln} in \inst{ex6_2_5}, \inst{ex6_2_7} and \inst{etamac};
\texttt{sin} and \texttt{cos} in \inst{lnts}, \inst{powerflow0030p} and \inst{powerflow0039p}, and \texttt{cos} also in \inst{hvycrash};
\texttt{tanh} in \inst{ann_cumene_tanh}.

\paragraph{GAMS and OSIL forms.}
For \cref{tab:sem-gams}: the GAMS files keep the objective variable that the OSIL files usually substitute out (\inst{dtoc5} has 100,000 GAMS and 99,999 OSIL variables).
The evaluations at random points used 50- or 60-digit arithmetic; for \inst{lnts50}--\inst{lnts400}, \inst{lukvle10}, \inst{chain50}--\inst{chain400}, \inst{ann_cumene_tanh} and the KAN instances, the GAMS files were first converted by GAMS, and their variables, types and bounds were compared exactly.
The exact comparisons checked every row, objective coefficient and bound of the GAMS text for \inst{dtoc5} and \inst{optcdeg2}; the constants and rows for \inst{camshape100}--\inst{camshape800}; the rows term by term for the three \inst{powerflow} models; the OSIL decoder against a separately written GAMS reader for the three \inst{eg} models; and, for \inst{waterno2_06}--\inst{waterno2_24}, the rows as polynomials together with the bounds, types and objective, a check that detected both of its negative controls.

\paragraph{Data reading by family.}
The binary64 reading (\cref{app:semantics-codes}) occurs only in early first-implementation scripts, where it is harmless: the \inst{lnts} certificate uses only binary64-exact data (its angle bound $\pm1.5707963267949$ is not used), \inst{lukvle10} has only integer data, \inst{dtoc5} enters $h=1/50000$ exactly, the \inst{camshape} strings are recovered exactly from their binary64 values by a shortest round trip, and the \inst{optcdeg2} script uses string literals.
The fast mode of the \inst{eg} search, a cross-check only (\cref{app:eg}), uses a rounded reading.
The number of distinct decimal strings that \reading{c} changes (\cref{app:semantics-binary64}) is 0 in \inst{lukvle10}, 2 in each \inst{lnts} and \inst{chain} file, between 3 and 10 in \inst{dtoc5}, \inst{optcdeg2}, \inst{catmix}, \inst{camshape} and \inst{hvycrash}, and from 37 (\inst{pindyck}) to 3,332 (\inst{eg_disc2_s}) in the other files.
The KAN and \inst{ann_cumene_tanh} files contain about 1,000 to 2,400 strings with 16 or 17 significant digits each, and \inst{powerflow0039r} contains 24; these strings are printed binary64 numbers.
\Cref{tab:trust-full} gives the data reading of each certificate family together with its trust base and primal construction.

\subsection{Claim register}\label{app:repro-register}

\Cref{tab:repro-register} maps each result to its checkers, recorded time, replay tier, evidence level, status and trust-base terms.
\nolinkurl{$P/artifact/RUNS.md} gives for every row of \cref{tab:repro-register} the full checker paths, arguments, expected outputs and run records, and \nolinkurl{$P/artifact/HISTORY.md} lists earlier and superseded certificates and displays, including strings in older records that must not be read as bounds.
Where an expected output reads ``as in \cref{tab:closures}'', the checker's exact output, rounded as stated in \cref{sec:semantics-display}, gives the table entry, and \nolinkurl{make_tables.py} checks this.
The input models are those of \cref{tab:sem-hashes}.

{\scriptsize
\setlength{\tabcolsep}{3pt}
\begin{longtable}{@{}>{\raggedright\arraybackslash}p{2.4cm}>{\raggedright\arraybackslash}p{7.6cm}>{\raggedright\arraybackslash}p{2.0cm}>{\raggedright\arraybackslash}p{3.3cm}@{}}
\caption{Claim register. ``2nd'': the other implementation and what it certifies; for \inst{optcdeg2}, \inst{lukvle10}, \inst{catmix}, \inst{etamac}, \inst{pindyck} and the \inst{emfl} enclosures it is the first implementation, for \inst{ex6_2_5}, \inst{ex6_2_7} and \inst{pricing050} a third code. Time: recorded wall time. Last column: tier (\cref{sec:repro}) of the replay of the displayed result by the checkers named first; evidence level; status (\cref{sec:semantics-protocol}; \emph{weaker second} and \emph{partial second} as in \cref{tab:trust}); trust-base terms besides \trust{read}, \trust{code} and \trust{thm}.}\label{tab:repro-register}\\
\toprule
result & checkers & time & tier; level; status; trust \\
\midrule
\endfirsthead
\multicolumn{4}{@{}l}{\cref{tab:repro-register}, continued}\\
\toprule
result & checkers & time & tier; level; status; trust \\
\midrule
\endhead
\bottomrule
\endfoot
\Cref{thm:lnts-opt}, \inst{lnts50}--\inst{lnts400} &
\nolinkurl{verify.py}; 2nd: \nolinkurl{lnts_exact_opt.py} (same optima) &
seconds &
1; \evid{hand}+\evid{stored}; verified; \trust{int} \\
\Cref{thm:dtoc5-bracket}, \inst{dtoc5} &
\nolinkurl{review.py}; 2nd: \nolinkurl{dtoc5_checks.py} (same display) &
17\,s; 3\,s &
1; \evid{hand}+\evid{stored}; verified; \trust{int} \\
\Cref{thm:optcdeg2-bound}, \inst{optcdeg2} &
\nolinkurl{v_states.py}, \nolinkurl{v_qcal_exact.py}; 2nd: \nolinkurl{optcdeg2_qcal_certify.py} (weaker). Primal: \nolinkurl{optcdeg2_primal_int.py}, \nolinkurl{v_primal.py} &
6\,s + 19\,s; 8\,s &
1; \evid{stored}; weaker second; \trust{int} (2nd primal: \trust{mp}) \\
\Cref{thm:lukvle10-bracket}, \inst{lukvle10} &
\nolinkurl{v_lukvle10_prep.py}, \nolinkurl{v_lukvle10_bnb.py}; 2nd: \nolinkurl{lukvle10_bound.py} (weaker). Primal: \nolinkurl{lukvle10_enclose.py}, \nolinkurl{lukvle10_crosscheck.py} &
5\,min; 6\,min &
1; \evid{rerun}; weaker second; \trust{mp} (dual), \trust{int} (primal) \\
\Cref{thm:chain-bound}, \inst{chain50}--\inst{chain400} &
\nolinkurl{chain_bound.py}; 2nd: \nolinkurl{v_chain_bnb.py} (same bound). Primal: \nolinkurl{build_points.py}, \nolinkurl{verify_points.py} &
12--29\,s; 72--195\,s &
1; \evid{rerun}; verified; \trust{mp} (dual), \trust{int} (primal) \\
\Cref{thm:catmix-bound}, \inst{catmix100}--\inst{catmix800} &
\nolinkurl{v_catmix_dp.py}, \nolinkurl{recheck_dp.py}; 2nd: \nolinkurl{catmix_bound.py} (weaker). Primal: \nolinkurl{exact_catmix.py}, \nolinkurl{policy_exact.py} &
20--46\,min; 17--77\,min &
2 (2nd: 3 for \inst{catmix800}); \evid{rerun} (dual), \evid{stored} (primal); weaker second; \trust{fp} \\
\Cref{thm:waterno2-bounds}, \inst{waterno2_06} &
\nolinkurl{verify_cs.py}; 2nd: \nolinkurl{ind_verify_cs.py}; pair bounds: \nolinkurl{vrebound_cs.py} &
13\,s; 11\,s; pairs about 60 CPU-h &
3 (path check: 1); \evid{rerun}; verified; \trust{fp}, \trust{int} \\
\Cref{thm:waterno2-bounds}, \inst{waterno2_09}--\inst{waterno2_24} &
\nolinkurl{run_period.py} (63 periods), \nolinkurl{vsum2.py}; regeneration by the first code: \nolinkurl{certify.py}. Primal: \nolinkurl{water_exact.py}, \nolinkurl{check_water_point.py} &
35--212\,min (7.4\,h in all) &
2 (09, 12), 3 (18, 24); \evid{rerun}; verified; \trust{fp}, \trust{int} \\
\Cref{thm:camshape-opt}, \inst{camshape100}--\inst{camshape800} &
\nolinkurl{v_camshape.py}; 2nd: \nolinkurl{check_exact.py}; third: \nolinkurl{mine.py} (directed rounding) &
8\,s; 6\,min; under 2\,s &
1; \evid{hand}+\evid{stored}; verified; \trust{int} \\
\Cref{thm:ex62-bounds}, \inst{ex6_2_5}, \inst{ex6_2_7} &
\nolinkurl{gibbs_kkt.py}, \nolinkurl{gibbs_bb.py}, \nolinkurl{gibbs_bound.py}; 2nd: \nolinkurl{gibbs_cert.py} (stronger). Primal: \nolinkurl{ex62_check.py} &
258\,s, 77\,s; 198\,s, 63\,s &
1; \evid{rerun}; verified; \trust{mp} (dual), \trust{int} (primal) \\
\Cref{thm:pricing-bound}, \inst{pricing050} &
\nolinkurl{v_pricing050.py}; 2nd: \nolinkurl{r2/pricing_check.py} (same bound). Primal: the earlier \nolinkurl{pricing_check.py} and \nolinkurl{r2/pricing_check.py} &
28\,s; 3\,s &
1; \evid{stored}; verified; \trust{mp} \\
\Cref{thm:etamac-bound}, \inst{etamac} &
\nolinkurl{v_etamac.py}; 2nd: \nolinkurl{etamac.py} (weaker). Primal: \nolinkurl{etamac_point.py} &
22\,s; 27\,s &
1; \evid{stored}; weaker second; \trust{mp} \\
\Cref{thm:pindyck-bound}, \inst{pindyck} &
\nolinkurl{own_ranges.py}, \nolinkurl{own_concavity.py}, \nolinkurl{verify_own_leaves.py}, \nolinkurl{final_bound.py}, \nolinkurl{primal_check.py}; 2nd: \nolinkurl{pindyck_global.py} (weaker); \nolinkurl{pindyck_sc.py} &
322\,s; 19\,s; under 1\,s &
1; \evid{stored}; weaker second; \trust{fp}, \trust{mp} \\
\Cref{thm:pf-0030p}, \inst{powerflow0030p} &
\nolinkurl{run_root.py}. Primal: \nolinkurl{certify.py}, \nolinkurl{verify.py} &
4\,s &
1; \evid{stored}; verified; \trust{int} \\
\Cref{thm:pf-0039}, \inst{powerflow0039p}, \inst{powerflow0039r} &
\nolinkurl{verify_exact.py}, \nolinkurl{verify_bb3.py}; 2nd: the leaf checks in \nolinkurl{powerflow0039-review-checks/} &
78\,s, 24\,s; 2\,s &
1; \evid{stored}; verified; \trust{int} \\
\Cref{prop:hvycrash-identity}, \inst{hvycrash} &
\nolinkurl{hvycrash.py}; 2nd: \nolinkurl{v_hvycrash.py}; \nolinkurl{hvy_check.py} &
under 1\,s each &
1; \evid{hand}; proved; \trust{int}, \trust{mp} (witness) \\
\Cref{thm:eg-bounds}, \inst{eg_int_s}, \inst{eg_disc_s}, \inst{eg_disc2_s} &
\nolinkurl{run_audit.py}, \nolinkurl{compare_audit.py} (leaf re-certification under the accuracy auditor); \nolinkurl{summarize.py}; coverage: \nolinkurl{own_cover.py}; search: \nolinkurl{egbb.py}. Primal: \nolinkurl{eg_dyadic_check.py}, \nolinkurl{exp_test.py}, \nolinkurl{verify_primal.py} &
16,703\,s over 44 jobs &
1, 2, 3 (by instance); \evid{rerun} (leaf partition stored); verified (\inst{eg_disc2_s}: partial second); \trust{fp}, accuracy auditor (\cref{thm:eg-trust}) \\
\Cref{thm:ann-bound}, \inst{ann_cumene_tanh} &
\nolinkurl{replay.py}, \nolinkurl{leaves.py}, \nolinkurl{verify_boxes.py}. Primal: \nolinkurl{nn_exact.py} &
33\,min, 1.6\,h; 1.2\,h (2 processes) &
3; \evid{rerun}; proved (the second code copied an idiom of the first); \trust{fp} \\
\Cref{prop:kan-infeasible}, six KAN instances &
\nolinkurl{kan_infeas_cert.py}; 2nd: \nolinkurl{kan_infeas_edge.py}, \nolinkurl{kan_infeas_edge_class.py} &
0.6--1.7\,s; about 1\,s &
1; \evid{hand}+\evid{stored}; verified; \trust{int} \\
\Cref{thm:kan-enclosure}, six KAN instances &
\nolinkurl{run_kan.py}; 2nd: the guarded \nolinkurl{kan_bnb_rigexp.py} through \nolinkurl{run_replays.py}; \nolinkurl{check_constants.py}. Points of $\RP$: \nolinkurl{nn_exact.py} &
22\,s--18\,min; 23\,s--20\,min &
1 (\texttt{r3}), 2 (\texttt{r5}); \evid{rerun}; verified (shared exponential and interval core); \trust{fp} \\
\Cref{sec:points}, \cref{tab:points} &
per family in the point folders \nolinkurl{lnts}, \nolinkurl{dtoc5-lukvle10}, \nolinkurl{chain}, \nolinkurl{powerflow}, \nolinkurl{water-ann-kan}, with separate checks; others as in the rows above &
seconds to minutes &
1; \evid{stored}; verified (\inst{etamac}: proved); \trust{int} (\trust{mp} for four points) \\
\Cref{prop:audit-refute}, \cref{tab:audit-pairs} &
\nolinkurl{audit.py}, \nolinkurl{verify_one.py}, \nolinkurl{cert_linear.py}, \nolinkurl{cert_ndnetgen.py}, \nolinkurl{cert_topopt.py}; 2nd: \nolinkurl{bound-audit-verification/}, \nolinkurl{bound-audit-recheck/} &
seconds per point; topopt about 3\,min &
1; \evid{hand}+\evid{stored} (\inst{topopt-cantilever_60x40_50}: \evid{rerun}); verified; \trust{int}; Krawczyk pairs also \trust{fp}, \trust{mp} \\
\Cref{prop:spring-opt,prop:emfl-enclosure}; \inst{rocket100}--\inst{rocket400} &
\nolinkurl{spring_global.py}; \inst{emfl}: \nolinkurl{emfl_bounds.py}, 2nd: \nolinkurl{cert_socp.py} (wider enclosures); \inst{rocket}: \nolinkurl{v_lindo.py}, 2nd: \nolinkurl{rocket_kraw.py} &
2--13\,s; 3--26\,s per \inst{emfl} instance &
1; \evid{hand}+\evid{stored} (\inst{emfl}: \evid{rerun}); verified (\inst{emfl} enclosures: weaker second); \trust{int} (\inst{rocket}: \cref{app:audit}) \\
\Cref{prop:scip-witnesses,prop:scip-reproducers,lem:scip-cube} &
\nolinkurl{exact_check.py}, \nolinkurl{spec_check.py}, \nolinkurl{indep_check.py}, \nolinkurl{gams_check.py}; SCIP's check: \nolinkurl{checksol_witness.py}; \nolinkurl{binary64_analysis.py}; reproducers \nolinkurl{fm336_v1010.cip}, \nolinkurl{tiny2.cip} &
seconds; SCIP runs need the listed builds &
1; \evid{stored}; verified; \trust{int} \\
\Cref{prop:baron-camshape,prop:qplib-copies,prop:kan-scip,prop:camino-eg,prop:minotaur-optcdeg2}, \cref{tab:claims} &
exact comparison in \nolinkurl{make_tables.py}; campaign points: \nolinkurl{point_checks.json} &
seconds; QPLIB copy optima under 2\,s and about 15\,min (two parsers) &
1; \evid{stored}; proved; \trust{int} \\
\Cref{tab:solvers,tab:claims-campaign} &
run logs, \nolinkurl{results_table.csv}; a rerun needs GAMS~54.3.1 and licences &
129 one-hour runs &
--; --; floating-point output; -- \\
Generated certified displays (\cref{app:displays}) &
\nolinkurl{make_tables.py} &
about 5\,s &
1; \evid{stored}; proved; \trust{int} \\
\end{longtable}
}

\subsection{Replay versus regeneration}\label{app:repro-regen}

\Cref{sec:repro} defines replay and regeneration; for \inst{lnts50}--\inst{lnts400}, \inst{dtoc5}, \inst{camshape} and \inst{hvycrash} there is no search, so replay is the whole check.

\paragraph{waterno2.}
The first code regenerates the period certificates: it derives each period's target from time-limited SCIP runs and certifies it with its branch and bound.
A clean-copy run reproduced the values and node counts of \inst{waterno2_06}, \inst{waterno2_09} and \inst{waterno2_12} exactly.
For \inst{waterno2_18} (periods $t=1$ and $t=7$) and \inst{waterno2_24} (period $t=14$), SCIP returned other incumbents, feasible within its tolerance of $10^{-6}$ and below the stored period bounds by about \sci{4.0}{-4}, \sci{3}{-5} and \sci{8.5}{-4}; the regenerated bounds, $4790.820086$ and $6576.150434$ (rounded down), are valid and slightly weaker than the reported $4790.820715$ and $6576.151388$.
Since both branch-and-bound codes certified the stored bounds, these incumbents are not exactly feasible (\cref{lem:sem-enclosure}(b)).
The reported values rest on the original run and on the separately written branch and bound with exact rational node bounds, both at the stored targets, and a clean-copy replay at these targets certified all 63 periods of \inst{waterno2_09}--\inst{waterno2_24} again.

\paragraph{optcdeg2 and powerflow.}
Regenerating the \inst{optcdeg2} split may give a different valid bound; the replay uses the stored split data.
Regeneration gives other multipliers and a weaker bound for \inst{powerflow0030p}, whose archived certificate replays exactly and gives the displayed value; a new SDP solve overwrites the stored certificate.
Regenerating the \inst{powerflow0039p} and \inst{powerflow0039r} leaf certificates (316\,s and 463\,s) gave identical final bounds.

\paragraph{ann\_cumene\_tanh.}
The certificate comes from two time-limited runs (1,800\,s and 5,400\,s).
Machine load changes the final frontier of a run stopped by a wall-clock limit, so the certificate is replayed by saved node counts, not regenerated.
The region lists of the second code's verification were not kept; the replay regenerates them, and the unsplit leaf boxes can be verified directly from the stored frontier.

\paragraph{eg.}
The audited run regenerated the recorded trees of \inst{eg_int_s} and \inst{eg_disc_s}, whose recording logs agree line by line with those of the original runs apart from timings; for \inst{eg_disc2_s} it re-certified the stored leaves, with per-leaf results identical byte for byte to those of the original run (\cref{app:egrounding-run}).
Coverage of every part is proved by an exact check that needs no tree.

\paragraph{KAN.}
A separate agent session reran the second code on the three \texttt{r3} models and reproduced every stored field except elapsed time; it checked the stored outputs of the three \texttt{r5} searches for consistency only.
A later replay of all six searches with the guarded quadratic bound routine reproduced every archived bound bit for bit, with no guard trigger (\cref{app:annkan-kan-enclosure}); the register row of \cref{thm:kan-enclosure} names its files, and the README gives the commands that stage and rerun this replay one instance at a time and compare each result with the archived one.

\paragraph{catmix, chain and the audit.}
The per-stage chord values of the \inst{catmix} certificates are not stored, so a check reruns the dynamic program (20 to 46 minutes per instance); the \inst{chain} end-value searches are short and rerun in full.
The \inst{topopt-cantilever_60x40_50} existence certificate of the audit does not store its basis, so it is regenerated, not replayed; the regenerated point also satisfies its row of \cref{tab:audit-pairs} (\cref{app:audit-existence}).
The \inst{emfl} certificates store the enclosures, not the rational dual vectors; a rerun checks new vectors exactly and may give slightly different, equally valid endpoints (\cref{app:audit-emfl}).
